\section*{第 \ref{ch:g12:satellites-kepler} 章——卫星与行星运动}

\begin{solution}{exo:g12:satellites-kepler:1}
在地面上：$F = GMm/R^2 = \num{3.98e17}/\num{4.06e13} \approx
\qty{9.8e3}{N}$。当 $r = \qty{6.77e6}{m}$ 时：
$F \approx \qty{8.7e3}{N}$——仍为地面值的 $89\%$。
航天员之所以漂浮，是因为他们与空间站一起\emph{下落}，
而不是因为\omterm{def:g10:universal-gravitation:weight}{重力}消失了。
\end{solution}

\begin{solution}{exo:g12:satellites-kepler:2}
$v = \sqrt{GM/R} = \sqrt{\num{3.98e14}/\num{6.37e6}} \approx
\qty{7.9e3}{m/s}$; $T = 2\pi R/v \approx \qty{5.1e3}{s} \approx
\qty{84}{min}$。这恰好是炮弹的 \qty{7.9}{km/s}：贴地
\omterm{def:g10:universal-gravitation:satellite}{轨道}\emph{正是}\omterm{def:g10:inertia:force}{牛顿}的那发炮弹。
\end{solution}

\begin{solution}{exo:g12:satellites-kepler:3}
$v \propto 1/\sqrt{r}$：除以 $\sqrt 2 \approx 1.41$。
$T \propto r^{3/2}$：乘以 $2^{3/2} \approx 2.8$。二者
均不含 $m$：否。
\end{solution}

\begin{solution}{exo:g12:satellites-kepler:4}
$v = \sqrt{\num{3.98e14}/\num{6.77e6}} \approx \qty{7.67e3}{m/s}$;
$T = 2\pi r/v \approx \qty{5.55e3}{s} \approx \qty{92}{min}$;
每天绕行 $\num{86400}/5546 \approx 15.6$ 圈。
\end{solution}

\begin{solution}{exo:g12:satellites-kepler:5}
\omterm{def:g10:universal-gravitation:satellite}{轨道}为以太阳为\omterm{def:g12:satellites-kepler:ellipse}{焦点}的\omterm{def:g12:satellites-kepler:ellipse}{椭圆}；相等时间内扫过相等面积；
$T^2/a^3$ 对所有行星都相同。速率在\omterm{def:g12:satellites-kepler:ellipse}{近日点}最大，由
开普勒第二定律（短弧段须扫得快）。否：\omterm{def:g10:universal-gravitation:satellite}{卫星}的
质量\omterm{def:g10:inertia:compensate}{相互抵消}——只有中心天体的质量起作用。
\end{solution}

\begin{solution}{exo:g12:satellites-kepler:6}
$T = \qty{2.36e6}{s}$: $M = 4\pi^2 r^3/(GT^2) =
\num{2.24e27}/371 \approx \qty{6.0e24}{kg}$——与
\qty{5.97e24}{kg} 相差在 $1\%$ 以内。
\end{solution}

\begin{solution}{exo:g12:satellites-kepler:7}
\omterm{def:g10:universal-gravitation:weight}{重力}指向地心，而圆周运动中\omterm{def:g11:forces-and-motion:netforce}{合力}指向
圆心：两个中心必须重合。位于巴黎纬线上方的
圆，其圆心在地轴上而非地心，这是不可能的。
悬停只有在纬线是以地心为圆心的大圆时才可行，
即在\emph{赤道}上空。
\end{solution}

\begin{solution}{exo:g12:satellites-kepler:8}
$T = \qty{2.754e4}{s}$: $M = 4\pi^2 r^3/(GT^2) =
\num{3.26e22}/\num{5.06e-2} \approx \qty{6.4e23}{kg}$——数据卡上的
火星。
\end{solution}

\begin{solution}{exo:g12:satellites-kepler:9}
地球：$T^2/a^3 = \num{9.96e14}/\num{3.35e33} =
\qty{2.98e-19}{s^2/m^3}$。火星：$T = \qty{5.94e7}{s}$，
$\num{3.52e15}/\num{1.18e34} = \qty{2.98e-19}{s^2/m^3}$——相等，符合
第三定律的要求。则 $M = 4\pi^2/(G \times \num{2.98e-19})
\approx \qty{1.99e30}{kg}$。
\end{solution}

\begin{solution}{exo:g12:satellites-kepler:10}
$v = 2\pi r/T = \qty{1.02e3}{m/s}$; $a = v^2/r =
\qty{2.72e-3}{m/s^2}$; $g/60^2 = 9.81/3600 = \qty{2.73e-3}{m/s^2}$.
月球位于 $60$ 个地球半径之外，其下落
加速度比苹果小 $60^2$ 倍：正是平方反比——同一定律
适用于两者。
\end{solution}

\begin{solution}{exo:g12:satellites-kepler:11}
谷神星：$T = 2.77^{3/2} = \sqrt{21.3} \approx \qty{4.6}{yr}$。哈雷彗星：
$T = 17.8^{3/2} \approx \qty{75}{yr}$，故下一次\omterm{def:g12:satellites-kepler:ellipse}{近日点}约在
$1986 + 75 \approx 2061$ 年。
\end{solution}

\begin{solution}{exo:g12:satellites-kepler:12}
在 \qty{400}{km} 处：$v \approx \qty{7.67e3}{m/s}$；在 \qty{350}{km}
（$r = \qty{6.72e6}{m}$）：$v \approx \qty{7.70e3}{m/s}$。\omterm{def:g10:inertia:inventory}{摩擦力}
确实消耗\omterm{def:g11:mechanical-energy:em}{机械能}——\omterm{def:g10:universal-gravitation:satellite}{轨道}缩小；但在下降途中
\omterm{def:g10:universal-gravitation:weight}{重力}做的功足以弥补阻力造成的速率损失：
$E_m$ 减小，而 $E_k$ 增大。
\end{solution}

\begin{solution}{exo:g12:satellites-kepler:13}
$T = \qty{3.024e5}{s}$: $M = 4\pi^2 r^3/(GT^2) =
\num{1.65e31}/6.10 \approx \qty{2.7e30}{kg} \approx 1.4$ 个太阳质量。
\emph{行星}的质量无法求出：它在\omterm{def:g10:universal-gravitation:satellite}{轨道}
方程中已被消去。
\end{solution}

\begin{solution}{exo:g12:satellites-kepler:14}
$T = \qty{8.862e4}{s}$, $GM = \qty{4.28e13}{m^3/s^2}$:
$r^3 = GMT^2/(4\pi^2) = \num{8.52e21}$，故 $r \approx
\qty{2.04e7}{m}$，且 $h = r - R \approx \qty{1.7e7}{m} \approx
\qty{17000}{km}$，即高出火星地面。
\end{solution}

\begin{solution}{exo:g12:satellites-kepler:15}
$v_a = r_p v_p / r_a = \num{8.77e10} \times \num{5.45e4} /
\num{5.25e12} \approx \qty{9.1e2}{m/s}$。由于 $r_a/r_p \approx 60$，
彗星的距离远 $60$ 倍，速率慢 $60$ 倍：它在
黑暗中徘徊数十年，却在数周内疾驰掠过太阳。
\end{solution}

\begin{solution}{pb:g12:satellites-kepler:1}
\textbf{1.} $T = \qty{86164}{s}$：地面随地球
\emph{相对于恒星}每\omterm{def:g12:satellites-kepler:geo}{恒星日}转一周；\qty{24}{h}
的一天多出的那一度，正是地球绕太阳每日前进的角度。

\textbf{2.} \omterm{def:g12:kinematics-2d:ucm}{匀速圆周运动}的\omterm{def:g11:forces-and-motion:netforce}{合力}指向圆
心；唯一的力——引力——指向地心，
两个圆心是同一个点。

\textbf{3.} 悬停于巴黎上空的圆周，其圆心应位于地心以北的
地轴上，由 2 可排除。
因此\omterm{def:g10:universal-gravitation:satellite}{轨道}必须位于\emph{赤道面}内。

\textbf{4.} \omterm{def:g10:inertia:force}{牛顿}第二定律的法向分量：
$m v^2/r = GmM/r^2$，故 $v = \sqrt{GM/r}$（$m$ 约去）。

\textbf{5.} $T = 2\pi r/v = 2\pi\sqrt{r^3/(GM)}$；两边平方，
$T^2/r^3 = 4\pi^2/(GM)$。

\textbf{6.} $r = \left(GMT^2/4\pi^2\right)^{1/3} =
(\num{7.49e22})^{1/3} \approx \qty{4.22e7}{m}$。

\textbf{7.} $h = r - R = \num{4.22e7} - \num{6.37e6} \approx
\qty{3.58e7}{m} \approx \qty{35800}{km}$。

\textbf{8.} $v = 2\pi \times \num{4.22e7}/\num{86164} \approx
\qty{3.07e3}{m/s}$。

\textbf{9.} $\sqrt{\num{3.98e14}/\num{4.22e7}} \approx
\qty{3.07e3}{m/s}$——一致。

\textbf{10.} 半径相同：质量已在 4 中约去。
质量较大的\omterm{def:g10:universal-gravitation:satellite}{卫星}只改变发射成本——
同一\omterm{def:g10:universal-gravitation:satellite}{轨道}需要更多燃料。

\textbf{11.} $v = \sqrt{\num{3.98e14}/\num{6.77e6}} \approx
\qty{7.67e3}{m/s}$。

\textbf{12.} $T = 2\pi r/v \approx \qty{5.55e3}{s} \approx
\qty{92.4}{min}$。

\textbf{13.} $\num{86400}/5546 \approx 15.6$ 圈\omterm{def:g10:universal-gravitation:satellite}{轨道}：
每天约 $16$ 次日出。

\textbf{14.} $r_{\text{地球静止}}/r_{\text{国际空间站}} = \num{4.22e7}/\num{6.77e6}
= 6.23$；$6.23^{3/2} \approx 15.5 = \num{86164}/5546$——开普勒第三定律，
此处即可验证。

\textbf{15.} 空间站的速率更大（\qty{7.67}{km/s}
对 \qty{3.07}{km/s}），但\omterm{def:g12:satellites-kepler:geo}{地球静止}\omterm{def:g10:universal-gravitation:satellite}{卫星}每\omterm{def:g10:orders-of-magnitude:unit}{千克}
需要更多\omterm{def:g10:energy-conservation:energy}{能量}：势能的增加超过了速率的减少。

\textbf{16.} $M_J = 4\pi^2 r^3/(G T^2)$。

\textbf{17.} $T = \qty{1.529e5}{s}$：$M_J =
\num{2.97e27}/1.56 \approx \qty{1.90e27}{kg}$。

\textbf{18.} $\num{1.90e27}/\num{5.97e24} \approx 318$ 个地球——
并约为太阳的 $1/1000$。

\textbf{19.} $T_J = \qty{3.576e4}{s}$：$r^3 = GM_J T_J^2/(4\pi^2) =
\num{4.11e24}$，$r \approx \qty{1.60e8}{m}$——木卫一位于
\qty{4.22e8}{m}，飞得高出 $2.6$ 倍，缓缓掠过木星
的天空。

\textbf{20.} 将中继站放置于高度为
\qty{35800}{km} 的赤道圆\omterm{def:g10:universal-gravitation:satellite}{轨道}上，它在此以 \qty{3.07}{km/s} 运行，
始终不偏离碟形天线的指向；同一定律，代入木卫一的公转\omterm{def:g12:kinematics-2d:ucm}{周期}，
便称出木星的质量为 $\qty{1.90e27}{kg}$——相当于 $318$ 个地球——不费吹灰之力。
\end{solution}
